\documentclass[11pt,a4paper]{article}
\usepackage[margin=2.5cm]{geometry}
\usepackage{amsmath,amssymb,amsthm}
\usepackage{enumitem}
\usepackage{hyperref}
\usepackage{url}
\usepackage{graphicx}
\usepackage{algorithm}
\usepackage{algpseudocode}
\usepackage{float}
\usepackage{tikz}
\usepackage{booktabs}
\usepackage{array}
\usepackage{float}
\usepackage{booktabs}
\usetikzlibrary{arrows.meta, positioning, shapes.geometric, calc, fit, backgrounds,automata}
\tikzset{
  state/.style={circle, draw, minimum size=1.2cm, inner sep=2pt, font=\small, thick},
  bad/.style={circle, draw=red!80!black, fill=red!20, minimum size=1.2cm, inner sep=2pt, font=\small, thick},
  arrow/.style={-Stealth, thick},
  every picture/.style={line width=0.8pt}
}
\newtheorem{theorem}{Theorem}[section]
\newtheorem{lemma}[theorem]{Lemma}
\newtheorem{proposition}[theorem]{Proposition}
\newtheorem{corollary}[theorem]{Corollary}
\newtheorem{question}[theorem]{Question}
\newtheorem{definition}[theorem]{Definition}
\newtheorem{remark}[theorem]{Remark}
\newtheorem{example}[theorem]{Example}
\newtheorem{observation}[theorem]{Observation}
\newtheorem{assumption}[theorem]{Assumption}
\newcommand{\D}{\mathcal{D}}
\newcommand{\CS}{\mathrm{CS}}
\newcommand{\cf}{\mathrm{cf}}
\newcommand{\lin}{\mathrm{lin}}
\newcommand{\Ccf}[1]{\mathcal C_{#1}^{\cf}}
\newcommand{\Clin}[1]{\mathcal C_{#1}^{\lin}}
\newcommand{\pref}{\mathrm{pref}}
\newcommand{\suf}{\mathrm{suf}}
\newcommand{\keywords}[1]{\smallskip\noindent{\bfseries keywords:} #1}
\usepackage{appendix}
\usepackage{titlesec}
\newcommand{\setupappendix}{
    \titleformat{\section}
    {\normalfont\Large\bfseries}
    {Appendix \thesection: }
    {0em}
    {}
}
\title{Distributional Learning of Context-Free Languages under Fixed Finite-Monoid Typing}
\author{Takayuki Kuriyama\\
Independent Researcher\\
\texttt{growup.kuriyama@gmail.com}}
\date{\today}
\begin{document}
\maketitle

\begin{abstract}
We study distributional learning of context-free languages under a fixed recognizable congruence $\sim_h$ given as the kernel of an explicit finite monoid homomorphism $h:\Sigma^*\to M$. For this fixed-$h$ setting, we develop a finite typed reconstruction theory for context-free $\sim_h$-substitutable languages. Starting from a reduced context-free grammar, we introduce a typed refinement that records both yield types and outer context types, show that the relevant structure is concentrated in a finite typed reconstruction basis, and prove that this basis is exposed by a finite observation set. Occurrences of the same nonterminal symbol may therefore have to be separated when their outer $h$-contexts differ.

We then prove exact reconstruction from positive data. From any finite sample $K\subseteq\Sigma^*$, we construct a canonical hypothesis grammar $\hat G(K)$, and we show that once $K$ contains the finite observation set associated with the target typed grammar, $\hat G(K)$ generates the target language exactly. Consequently, for every explicit finite monoid homomorphism $h$, the class $\mathcal C_h^{\mathrm{cf}}$ of context-free $\sim_h$-substitutable languages is identifiable in the limit from positive data, with polynomial-time hypothesis construction and update.

For the linear subclass $\mathcal C_h^{\mathrm{lin}}$, we further prove polynomial upper bounds on characteristic-sample size and word length. Thus the same learner gives a full polynomial time-and-data result for linear targets.

\keywords{grammatical inference, distributional learning, context-free languages, recognizable congruences, finite monoid homomorphisms, identification in the limit, exact reconstruction}
\end{abstract}

\section{Introduction}
\label{sec:intro}

\subsection{Background and motivation}

Gold~\cite{gold1967} showed that no language class containing all finite
languages and at least one infinite language is identifiable in the limit
from positive data alone. This negative result motivated a broad line of
research in grammatical inference; see, for example,
de~la~Higuera~\cite{higuera2010} and
Heinz, de~la~Higuera, and van~Zaanen~\cite{heinz-higuera-zaanen2011}
for general background. For context-free languages, one standard way to
recover positive learning results is the \emph{distributional} approach,
in which learnability is obtained by restricting which substrings may be
compared across contexts. Classical examples include Clark--Eyraud
substitutability~\cite{clark-eyraud2007} and Yoshinaka's
$(k,\ell)$-substitutability~\cite{yoshinaka2008}. More generally,
positive-data learnability has also been established for other grammar
classes under suitable finite control assumptions; see, for example,
Kanazawa~\cite{kanazawa1998} for restricted categorial grammars.

The present paper studies a recognizable-congruence version of this
distributional program. Let $h:\Sigma^*\to M$ be a homomorphism into a
finite monoid, and write $\sim_h:=\ker h$. A language $L\subseteq\Sigma^*$
is \emph{$\sim_h$-substitutable} if $h(x)=h(y)$ and
$\D_L(x)\cap\D_L(y)\neq\emptyset$ imply $\D_L(x)=\D_L(y)$ for all
$x,y\in\Sigma^*$. Thus full equality of distributions is required only
for strings that already agree under a fixed finite algebraic summary.

Recognizable congruences are a natural comparison mechanism here. If
$L=h^{-1}(F)$ for some finite monoid homomorphism $h:\Sigma^*\to M$ and
some $F\subseteq M$, then any two strings with the same $h$-value
automatically have the same continuation behavior, and hence the same
distribution. In particular, every regular language belongs to the broader
class $\mathsf{RS}$ of languages that are $\sim_h$-substitutable for
\emph{some} finite monoid homomorphism. This observation shows that finite
monoid typings are not an ad hoc generalization of bounded contexts, but a
canonical algebraic source of finite distributional control.

Our focus is the \emph{fixed-$h$} regime: an explicit finite monoid
homomorphism $h:\Sigma^*\to M$ is given in advance, and the learner is
required to learn under that fixed comparison mechanism. In other words,
we separate two problems: determining from data which finite congruence
should govern comparison, and learning once that congruence has already
been specified. This separation is natural. Clark--Eyraud substitutability
and Yoshinaka's $(k,\ell)$-substitutability also fix in advance a finite
comparison mechanism; the present framework replaces bounded context
windows by an arbitrary recognizable congruence.

\paragraph{Relation to prior learners.}
At the level of surface rule schemata, the learner constructed here still
belongs to the distributional tradition: it builds a grammar from local
factorizations $uxv\in K$ extracted from a finite positive sample $K$.
The novelty is not this surface format, but the fixed-$h$ structural
theory behind it. An explicit homomorphism $h:\Sigma^*\to M$ induces a
typed refinement, a finite reconstruction basis, and a finite observation
set sufficient for exact recovery. Clark--Eyraud substitutability appears
as the trivial-monoid case. In the genuinely nontrivial fixed-$h$ setting,
by contrast, one must simultaneously control yield types and external
context types, and the correctness proof proceeds via typed transport and
typed substitution rather than the untyped Clark--Eyraud argument.


\subsection{Contributions and main results}

We do not address the fully adaptive problem of inferring a suitable
congruence from positive data alone. Instead, we isolate the learning
problem that remains once an explicit finite monoid homomorphism
$h:\Sigma^*\to M$ has been fixed. The resulting theory has two distinct
layers. For the general context-free class $\Ccf{h}$, we prove exact
reconstruction, identification in the limit from positive data, and
polynomial-time hypothesis construction and update, but we do \emph{not}
claim a polynomial upper bound on the size of a characteristic sample.
For the linear subclass $\Clin{h}$, we additionally prove polynomial upper
bounds on characteristic-sample size and word length, yielding a full
polynomial time-and-data result.

More precisely, the contributions of the paper are as follows.

\begin{enumerate}[label=\textup{(\arabic*)},leftmargin=2.4em]
\item We introduce a \emph{typed refinement} $\widetilde G$ of a reduced
context-free grammar under the fixed homomorphism $h:\Sigma^*\to M$.
This refinement separates occurrences of the same untyped nonterminal
according to both the $h$-type of the generated yield and the $h$-types
of the surrounding left and right contexts.

\item We define a \emph{finite typed reconstruction basis}
$\mathcal B(\widetilde G)$ consisting of the reachable productive typed
nonterminals, the realized typed rule instances, and the associated
canonical terminal and context data. This basis isolates the finite typed
local information needed for exact recovery of the target language.

\item We show that the typed reconstruction basis is \emph{exposed by a
finite observation set} $\CS(\widetilde G)$. Thus the hidden typed
structure is not merely postulated abstractly: it is forced by finitely
many observable words of the target language.

\item From a finite positive sample $K$, we construct a \emph{canonical
hypothesis grammar} $\hat G(K)$ and prove an \emph{exact reconstruction
theorem}: once $K$ contains $\CS(\widetilde G)$, the grammar $\hat G(K)$
generates the target language exactly.

\item As a learning-theoretic consequence, for every explicit finite
monoid homomorphism $h:\Sigma^*\to M$, the class $\Ccf{h}$ is
\emph{identifiable in the limit from positive data}. Moreover, for every
finite sample, the current hypothesis can be constructed and updated in
polynomial time in the current sample size and total input length.

\item For the linear subclass $\Clin{h}$, we prove polynomial upper
bounds on the size of a characteristic sample and on the lengths of its
words. Hence the same learner achieves a full \emph{polynomial
time-and-data} result on linear targets.
\end{enumerate}

We further situate the fixed-$h$ framework within the broader landscape of distributional learning. Writing $\mathsf{KL}$ for the union of Yoshinaka's fixed $(k,\ell)$-substitutable classes, we show that the capped counter family $\mathrm{CCL}_p$ already yields the strict inclusion $\mathsf{KL}\subsetneq\mathsf{RS}$ at the regular level. By contrast, several natural deterministic context-free languages lie outside $\mathsf{RS}$; examples include the uncapped counter language $\mathrm{CCL}$, the one-bracket Dyck language $D_1$, and Yoshinaka's classical language $L(S\to aSS\mid b)$.

We also give a positive answer to the question of whether $\mathsf{RS}\setminus\mathsf{KL}$ contains any non-regular context-free language. Concretely, we show that $L^*=\{a^n b^n:n\geq 0\}^*$ is a non-regular deterministic context-free language in $\mathsf{RS}\setminus\mathsf{KL}$, thereby establishing that the strict inclusion $\mathsf{KL}\subsetneq\mathsf{RS}$ extends beyond the regular level to the non-regular context-free level. These boundary results are secondary to the main positive results, but they clarify the scope of recognizable-congruence control. From this perspective, a structural characterization of $\mathsf{RS}\cap\mathrm{CFL}$ remains a natural problem for future work.

Taken together, these results show that the fixed-$h$ setting carves out a mathematically robust and structurally transparent subtheory within distributional context-free learning.

\subsection{Organization of the paper}
Section~\ref{sec:prelim} recalls the basic notions and positions the
classical substitutability frameworks inside the recognizable-congruence
setting. Section~\ref{sec:framework} formulates the fixed-$h$ model and
states the main theorem. Section~\ref{sec:typing} develops the typed
refinement and the finite typed reconstruction basis, and
Section~\ref{sec:algo} proves exact reconstruction and identification in
the limit. Section~\ref{sec:complexity} establishes polynomial-time upper
bounds for hypothesis construction and update, while Section~\ref{sec:linear}
proves polynomial time and data for the linear subclass.
Section~\ref{sec:ccl} and Section~\ref{sec:non-sub} provide the
structural separation and boundary results. Section~\ref{sec:lstar} resolves
the open question left by those boundary results by showing that
$L^*=\{a^nb^n:n\geq 0\}^*$ is a non-regular deterministic context-free
language in $\mathsf{RS}\setminus\mathsf{KL}$, establishing that the strict
inclusion $\mathsf{KL}\subsetneq\mathsf{RS}$ extends beyond the regular
level. Appendix~\ref{sec:properties} records additional closure properties
and counterexamples, and Appendix~\ref{app:ccl1} treats the exceptional
case of $\mathrm{CCL}_1$.

\section{Preliminaries}
\label{sec:prelim}

Let $\Sigma^*$ denote the free monoid over a finite alphabet $\Sigma$,
and let $\lambda$ denote its identity element. For a word
$w\in\Sigma^*$, we write $|w|$ for its length. For a language
$L\subseteq\Sigma^*$ and $x\in\Sigma^*$, the \emph{distribution} of $x$
in $L$ is defined by
$\D_L(x):=\{(u,v)\in\Sigma^*\times\Sigma^*:uxv\in L\}$.
We use standard asymptotic notation, and all claims about polynomial
time are understood with respect to the natural finite input size
specified locally at the relevant point.

\subsection{Learning from positive data}

Throughout the paper, we work in Gold's model of identification in the
limit~\cite{gold1967}. A \emph{text} for a language $L\subseteq\Sigma^*$
is an infinite sequence $T=(w_1,w_2,\dots)$ of words from $L$ such that
every word in $L$ appears at least once in $T$. For $n\ge 1$, we write
$T[n]:=(w_1,\dots,w_n)$ for the prefix of $T$ of length $n$.

A \emph{learner} is a computable mapping that assigns to each finite
text prefix a hypothesis grammar, or equivalently the language generated
by that grammar. A learner $\mathcal A$ \emph{identifies in the limit
from positive data} a language class $\mathcal C$ if for every
$L\in\mathcal C$ and every text $T$ for $L$, there exists $n_0$ such
that for all $n\ge n_0$, we have $L(\mathcal A(T[n]))=L$.

Since the learners considered in this paper produce their output
independently of the order in which positive examples are observed, we
freely identify a finite text prefix $(w_1,\dots,w_n)$ with the
corresponding finite sample $K=\{w_1,\dots,w_n\}\subseteq\Sigma^*$.

\subsection{Characteristic samples}

We adopt the standard characteristic-sample viewpoint in grammatical
inference~\cite{higuera1997,higuera2010}. Let $\mathcal C$ be a language
class and $\mathcal A$ a learner.

\begin{definition}[Characteristic sample]
Let $L\in\mathcal C$. A finite set $S\subseteq L$ is called a
\emph{characteristic sample} of $L$ with respect to $\mathcal A$ if
$L(\mathcal A(K))=L$ holds for every finite sample $K$ satisfying
$S\subseteq K\subseteq L$.
\end{definition}

\begin{definition}[Polynomial time-and-data, in the standard sense of de la Higuera]
Let $\mathcal{R}$ be a representation class for a language class $\mathcal{C}$, and let $L(r)$ denote the language represented by $r \in \mathcal{R}$. We say that a learner $\mathcal{A}$ \emph{identifies $\mathcal{R}$ from positive data in polynomial time and data} if there exist polynomials $p_{\mathcal{A}}$ and $q_{\mathcal{A}}$ such that the following hold:

(i) for every finite positive sample $K$, the hypothesis $\mathcal{A}(K)$ is computable in time at most $p_{\mathcal{A}}(\|K\|)$, where $\|K\| := \sum_{w \in K} |w|$;

(ii) for every target representation $r \in \mathcal{R}$, there exists a finite set $\CS(r) \subseteq L(r)$ such that $\|\CS(r)\| \le q_{\mathcal{A}}(|r|)$, and for every finite sample $K$ satisfying $\CS(r) \subseteq K \subseteq L(r)$, the hypothesis $\mathcal{A}(K)$ exactly reconstructs $L(r)$.
\end{definition}

\subsection{Recognizable congruences and fixed homomorphisms}

A congruence on $\Sigma^*$ is an equivalence relation compatible with
concatenation. Such a congruence is called \emph{recognizable} if it has
finite index, and this is equivalent to being the kernel of a monoid
homomorphism into a finite monoid~\cite{eilenberg1974}.

Let $M$ be a finite monoid and let $h:\Sigma^*\to M$ be a monoid
homomorphism. We write
$\sim_h:=\ker h=\{(x,y)\in\Sigma^*\times\Sigma^*:h(x)=h(y)\}$.

Throughout the paper, we assume that $h$ is given \emph{explicitly} to
the learner designer: the multiplication table of $M$ and the values
$h(a)\in M$ for all $a\in\Sigma$ are given as fixed prior data. Since
$M$ is fixed in the main theorem, the time required to compute $h(w)$
for $w\in\Sigma^*$ is linear in $|w|$.

\begin{definition}[\texorpdfstring{$\sim_h$}{sim\_h}-substitutability]
A language $L\subseteq\Sigma^*$ is said to be
\emph{$\sim_h$-substitutable} if for all $x,y\in\Sigma^*$,
\[
h(x)=h(y)\ \text{and}\ \D_L(x)\cap\D_L(y)\neq\emptyset
\Longrightarrow \D_L(x)=\D_L(y)
\]
holds.
\end{definition}

As noted above, we write $\Ccf{h}$ for the class of context-free
$\sim_h$-substitutable languages, and $\Clin{h}$ for its linear
subclass.

\section{The Framework of Fixed Finite-Monoid Typing}
\label{sec:framework}

Fix a finite monoid $M$ and a homomorphism $h:\Sigma^*\to M$. The aim
of this section is to place the fixed-$h$ setting in its proper
structural context. We first explain why recognizable congruences
provide a natural finite control on distributions, then show that the
classical frameworks of Clark--Eyraud and Yoshinaka arise as special
cases, and finally state the main fixed-$h$ theorem used throughout the
paper.

\subsection{Recognizable congruences as distributional control}
\label{subsec:recognizable-control}

Every homomorphism $h:\Sigma^*\to M$ induces an equivalence relation
$\sim_h$ on $\Sigma^*$ by $x\sim_h y \iff h(x)=h(y)$. Since
$h(uxv)=h(u)h(x)h(v)$ holds for all $u,x,v\in\Sigma^*$, the relation
$\sim_h$ is a monoid congruence of finite index. Such congruences are
precisely the recognizable congruences: conversely, every recognizable
congruence arises as the kernel of a homomorphism into a finite
monoid~\cite{eilenberg1974}.

The fixed-$h$ setting used below assumes that the homomorphism is given
\emph{explicitly}. Concretely, the multiplication table of $M$ and the
values $h(a)\in M$ for all $a\in\Sigma$ are included as part of the
prior data. Thus the learner is not required to infer the algebraic
typing from examples, but instead receives its finite control mechanism
in operational form.

Recall that $L\subseteq\Sigma^*$ is $\sim_h$-substitutable if
$h(x)=h(y)$ and $\D_L(x)\cap\D_L(y)\neq\emptyset$ together imply
$\D_L(x)=\D_L(y)$ for all $x,y\in\Sigma^*$.

In other words, two strings are required to have exactly the same full
distribution only when they are already equivalent under the fixed
finite algebraic summary given by $h$ and already share at least one
common context. This is the basic principle of distributional control in
the present paper.

For later reference, define
\[\mathsf{RS}:=\{L\subseteq\Sigma^*\mid\text{$L$ is $\sim_h$-substitutable
for some finite monoid homomorphism $h$}\},\]
the class of relation-substitutable languages.

\begin{proposition}\label{prop:regular-auto}
Let $h:\Sigma^*\to M$ be a finite monoid homomorphism, and let
$F\subseteq M$. If $L=h^{-1}(F)$, then $L$ is
$\sim_h$-substitutable. In particular, every regular language belongs to
$\mathsf{RS}$.
\end{proposition}

\begin{proof}
Let $x,y\in\Sigma^*$ satisfy $h(x)=h(y)$ and
$\D_L(x)\cap\D_L(y)\neq\emptyset$. Take any $(u,v)\in\D_L(x)$. Then
$uxv\in L=h^{-1}(F)$, so $h(u)h(x)h(v)\in F$. Since $h(x)=h(y)$, we
also have $h(u)h(y)h(v)\in F$, hence $uyv\in L$ and
$(u,v)\in\D_L(y)$. Therefore $\D_L(x)\subseteq\D_L(y)$, and the reverse
inclusion follows symmetrically.
\end{proof}

\subsection{Classical substitutability as special cases}
\label{subsec:classical-subst}

Clark--Eyraud substitutability corresponds to the trivial monoid case,
and each fixed $(k,\ell)$-substitutable class embeds into the present
framework via a canonical finite monoid quotient. We record this for
completeness.

\begin{definition}[Clark--Eyraud substitutability]
A language $L\subseteq\Sigma^*$ is \emph{substitutable} if for all
$x,y\in\Sigma^*$, $\D_L(x)\cap\D_L(y)\neq\emptyset$ implies
$\D_L(x)=\D_L(y)$.
\end{definition}

\begin{definition}[Yoshinaka's \texorpdfstring{$(k,\ell)$}{(k,l)}-substitutability]
A language $L\subseteq\Sigma^*$ is \emph{$(k,\ell)$-substitutable} if
for all $x_1,y_1,z_1,x_2,y_2,z_2\in\Sigma^*$, all $v\in\Sigma^k$, and
all $u\in\Sigma^\ell$, the conditions $vy_1u,vy_2u\neq\lambda$,
$x_1vy_1uz_1\in L$, $x_1vy_2uz_1\in L$, and $x_2vy_1uz_2\in L$ together
imply $x_2vy_2uz_2\in L$.
\end{definition}

In other words, $(k,\ell)$-substitutability requires a bounded-boundary
replacement principle: if two nonempty substrings $y_1$ and $y_2$ are both
admissible inside one common boundary pair $v\in\Sigma^k$, $u\in\Sigma^\ell$,
and if $y_1$ is admissible in another ambient context, then $y_2$ must also be
admissible in that second context\footnote{In particular, Yoshinaka's
$(0,0)$-substitutability may be written as: if $x,y\neq\lambda$ and
$\D_L(x)\cap\D_L(y)\neq\emptyset$, then $\D_L(x)=\D_L(y)$. Thus the
quantification is restricted to nonempty strings, so this condition is weaker
than Clark--Eyraud substitutability, which requires
$\D_L(x)\cap\D_L(y)\neq\emptyset \Rightarrow \D_L(x)=\D_L(y)$ for all
$x,y\in\Sigma^*$.}. We write
$\mathsf{KL}:=\bigcup_{k,\ell\ge0}\{L\subseteq\Sigma^*\mid\text{$L$ is
$(k,\ell)$-substitutable}\}$.

\begin{proposition}\label{prop:ce-special}
If $L$ is Clark--Eyraud substitutable, then there exists a finite monoid
homomorphism $h:\Sigma^*\to M$ such that $L$ is
$\sim_h$-substitutable. In particular, every Clark--Eyraud
substitutable language belongs to $\mathsf{RS}$.
\end{proposition}

\begin{proof}
Take the trivial monoid $M=\{1\}$ and let $h:\Sigma^*\to M$ be the
unique homomorphism. Then $h(x)=h(y)$ holds for all
$x,y\in\Sigma^*$, so $\sim_h$-substitutability reduces exactly to
Clark--Eyraud substitutability.
\end{proof}

\begin{proposition}\label{prop:yl-special}
For each fixed $k,\ell\ge0$, there exists a finite monoid homomorphism
$h_{k,\ell}:\Sigma^*\to M_{k,\ell}$ such that every
$(k,\ell)$-substitutable language is
$\sim_{h_{k,\ell}}$-substitutable. In particular, each fixed
$(k,\ell)$-substitutable class is contained in $\mathsf{RS}$.
\end{proposition}

\begin{proof}
Fix $k,\ell\ge0$, and write $N:=\max\{1,k+\ell\}$. Also adopt the
convention $\pref_0(w)=\suf_0(w)=\lambda$. Define
\[
T:=\Sigma^{<N}\sqcup(\Sigma^k\times\Sigma^\ell),
\]
and let $\tau:\Sigma^*\to T$ be given by
\[
\tau(w):=\begin{cases}
w & \text{if $|w|<N$,}\\
(\pref_k(w),\suf_\ell(w)) & \text{if $|w|\ge N$.}
\end{cases}
\]
Since $\Sigma$ is finite and $k,\ell$ are fixed, the set $T$ is finite.

Define a relation $\equiv$ on $\Sigma^*$ by
$x\equiv y \iff \tau(x)=\tau(y)$. We first show that $\equiv$ is a
congruence. Suppose that $\tau(x)=\tau(x')$ and $\tau(y)=\tau(y')$.
It suffices to prove that $\tau(xy)=\tau(x'y')$.

First consider the case $|xy|<N$. Then $|x|,|y|<N$, so
$\tau(x)=x$ and $\tau(y)=y$. Since $T$ is defined as the disjoint union
$\Sigma^{<N}\sqcup(\Sigma^k\times\Sigma^\ell)$, the equality
$\tau(x)=\tau(x')$ implies $x'=x$, and similarly $\tau(y)=\tau(y')$
implies $y'=y$. Hence $\tau(xy)=\tau(x'y')$.

We may therefore assume that $|xy|\ge N$. In this case $\tau(xy)$ is
determined by $(\pref_k(xy),\suf_\ell(xy))$, so it is enough to show
that both $\pref_k(xy)$ and $\suf_\ell(xy)$ are determined solely by
$\tau(x)$ and $\tau(y)$.

First consider the prefix. If $|x|\ge k$, then $\pref_k(xy)=\pref_k(x)$,
which is determined by $\tau(x)$. If $|x|<k$, then
$\pref_k(xy)=x\,\pref_{k-|x|}(y)$. In this case $|x|<k\le k+\ell\le N$,
so $\tau(x)=x$. Moreover, if $|y|<N$, then $\tau(y)=y$, and hence
$\pref_{k-|x|}(y)$ is clearly determined by $\tau(y)$. If $|y|\ge N$,
then $\tau(y)=(\pref_k(y),\suf_\ell(y))$, and since $k-|x|\le k$, the
word $\pref_{k-|x|}(y)$ is determined by $\pref_k(y)$, hence by
$\tau(y)$. Thus $\pref_k(xy)$ is always determined by $\tau(x)$ and
$\tau(y)$.

Next consider the suffix. If $|y|\ge\ell$, then $\suf_\ell(xy)=\suf_\ell(y)$,
which is determined by $\tau(y)$. If $|y|<\ell$, then
$\suf_\ell(xy)=\suf_{\ell-|y|}(x)\,y$. In this case
$|y|<\ell\le k+\ell\le N$, so $\tau(y)=y$. Moreover, if $|x|<N$, then
$\tau(x)=x$, so $\suf_{\ell-|y|}(x)$ is determined by $\tau(x)$. If
$|x|\ge N$, then $\tau(x)=(\pref_k(x),\suf_\ell(x))$, and since
$\ell-|y|\le\ell$, the word $\suf_{\ell-|y|}(x)$ is determined by
$\suf_\ell(x)$, hence by $\tau(x)$. Therefore $\suf_\ell(xy)$ is also
determined by $\tau(x)$ and $\tau(y)$.

It follows that $\tau(xy)$ is determined solely by $\tau(x)$ and
$\tau(y)$. Hence from $\tau(x)=\tau(x')$ and $\tau(y)=\tau(y')$ we
obtain $\tau(xy)=\tau(x'y')$. Thus $\equiv$ is a finite-index
congruence. Therefore the quotient
$M_{k,\ell}:=\Sigma^*/{\equiv}$ is a finite monoid, and the quotient map
$h_{k,\ell}:\Sigma^*\to M_{k,\ell}$ is a finite monoid homomorphism.
By construction, $h_{k,\ell}(s)=h_{k,\ell}(t)$ is equivalent to
$\tau(s)=\tau(t)$.

Now let $L\subseteq\Sigma^*$ be $(k,\ell)$-substitutable, and suppose
that $h_{k,\ell}(s)=h_{k,\ell}(t)$ and
$\D_L(s)\cap\D_L(t)\neq\emptyset$. We show that
$\D_L(s)=\D_L(t)$.

First consider the case $|s|<N$. Then $\tau(s)=s$, which lies in the
first component of the disjoint union $T$. Since $\tau(s)=\tau(t)$ and
$T$ is a disjoint union, it follows that $|t|<N$ and $t=s$. Hence
$\D_L(s)=\D_L(t)$ trivially.

We may therefore assume that $|s|,|t|\ge N$. Since $\tau(s)=\tau(t)$,
we have $\pref_k(s)=\pref_k(t)=:v$ and $\suf_\ell(s)=\suf_\ell(t)=:u$.
Thus we may write $s=vy_1u$ and $t=vy_2u$. Moreover,
$|s|,|t|\ge N\ge1$, so $vy_1u=s\neq\lambda$ and $vy_2u=t\neq\lambda$.

Take $(x_1,z_1)\in\D_L(s)\cap\D_L(t)$. Then
$x_1vy_1uz_1\in L$ and $x_1vy_2uz_1\in L$. Now let
$(x_2,z_2)\in\D_L(s)$ be arbitrary. Then $x_2vy_1uz_2\in L$. Applying
the definition of $(k,\ell)$-substitutability to
$x_1,y_1,z_1,x_2,y_2,z_2,v,u$, we obtain $x_2vy_2uz_2\in L$. Hence
$(x_2,z_2)\in\D_L(t)$. Therefore $\D_L(s)\subseteq\D_L(t)$.
The reverse inclusion follows symmetrically by exchanging $s$ and $t$.
Thus $\D_L(s)=\D_L(t)$.

Therefore $L$ is $\sim_{h_{k,\ell}}$-substitutable. In particular, the
fixed $(k,\ell)$-substitutable class is contained in $\mathsf{RS}$.
\end{proof}

\begin{corollary}\label{cor:framework-special-cases}
Clark--Eyraud substitutability and each fixed
$(k,\ell)$-substitutable class are special cases of the framework of
fixed recognizable congruences.
\end{corollary}

\begin{proof}
This follows by combining Proposition~\ref{prop:ce-special} and
Proposition~\ref{prop:yl-special}.
\end{proof}

\subsection{The main fixed-\texorpdfstring{$h$}{h} theorem}
\label{subsec:main-fixed-h}

Throughout the paper, we work under the fixed-$h$ framework: a specific
finite monoid homomorphism $h:\Sigma^*\to M$ is given explicitly as part
of the prior data and remains fixed throughout the entire learning
process. Equivalently, the learner is given the finite algebraic typing
in operational form and is not required to infer it from data.

Fix a finite monoid homomorphism $h:\Sigma^*\to M$. As recalled above,
we write
\[
\begin{aligned}
\Ccf{h}
&:=\bigl\{
L\subseteq\Sigma^*
\;\bigm|\;
\text{$L$ is context-free and $\sim_h$-substitutable}
\bigr\},\\
\Clin{h}
&:=\bigl\{
L\subseteq\Sigma^*
\;\bigm|\;
\text{$L$ is linear context-free and $\sim_h$-substitutable}
\bigr\}.
\end{aligned}
\]
Clearly, $\Clin{h}\subseteq\Ccf{h}$. The learnability of these classes
is the content of our main theorem below.

\begin{theorem}[Fixed-\texorpdfstring{$h$}{h} learning theorem]
\label{thm:main}
For every explicit finite monoid homomorphism $h:\Sigma^*\to M$, there
exists a learner $\mathcal A_h$ such that: \textup{(i)} $\mathcal A_h$
identifies in the limit from positive data every language in
$\Ccf{h}$; \textup{(ii)} for every finite sample, the hypothesis
constructed by $\mathcal A_h$ can be computed and updated in polynomial
time in the current sample size and total input length.
\end{theorem}

The proof for fixed $h$ is organized as a finite typed reconstruction
theory. Section~\ref{sec:typing} develops the typed local structure and
the associated finite reconstruction data, Section~\ref{sec:algo}
proves complete reconstruction from finite observations and derives
identification in the limit, and Section~\ref{sec:complexity}
establishes polynomial-time upper bounds for hypothesis construction and
update.

\begin{corollary}\label{cor:ce-yoshinaka}
Theorem~\ref{thm:main} specializes both to Clark--Eyraud substitutable
languages and to Yoshinaka's $(k,\ell)$-substitutable languages for each
fixed $k,\ell$.
\end{corollary}

\begin{proof}
This follows by combining Theorem~\ref{thm:main} with
Proposition~\ref{prop:ce-special} and
Proposition~\ref{prop:yl-special}.
\end{proof}

\begin{remark}\label{rem:main-poly-caveat}
Theorem~\ref{thm:main}\textup{(ii)} does \emph{not} claim a polynomial
upper bound on the size of a characteristic sample. Intuitively, beyond
the linear subclass, shortest distinguishing contexts may depend on
branching derivations or nested structural dependencies, so there is no
a priori reason to expect uniformly small samples. By contrast, when one
restricts to the linear subclass $\Clin{h}$, the constrained shape of
derivations makes it possible to prove that both the size of a
characteristic sample and the lengths of its words are bounded by a
polynomial in the grammar size (Theorem~\ref{thm:linear-poly}).
\end{remark}

The notation specific to typed reconstruction is summarized at the
beginning of Section~\ref{sec:typing}.

\section{Finite Typed Reconstruction under a Fixed Recognizable Congruence}
\label{sec:typing}

Under a fixed recognizable congruence, the relevant syntactic
information of a $\sim_h$-substitutable language can be compressed into
finitely many typed local configurations. The goal of this section is to
make this finite control explicit and to show that complete recovery of
the target language follows by reconstructing this finite typed
structure from positive data.

The notation in Table~\ref{tab:notation-typed} will be used throughout the remainder of this paper.

\begin{table}[h]
\centering
\small
\begin{tabular}{>{\raggedright\arraybackslash}p{0.22\linewidth} >{\raggedright\arraybackslash}p{0.70\linewidth}}
\toprule
Notation & Meaning \\
\midrule
$G=(V,\Sigma,P,S_0)$ & a reduced SSBNF grammar for the target language $L$ (Definition~\ref{def:SSBNF}) \\
$\widetilde G$ & the trimmed typed refinement of $G$ under the fixed homomorphism $h$ (Section~\ref{subsec:typed-refinement}) \\
$A_p^{m,n}$ & the typed copy of the nonterminal $A$, with yield $h$-type $p$ and surrounding-context $h$-type $(m,n)$ \\
$\omega(X)$ & the lexicographically least terminal yield derivable from the productive symbol $X$ \\
$\chi(X)=\langle u_X,v_X\rangle$ & the lexicographically least context pair satisfying $\widetilde S\Rightarrow_{\widetilde G}^* u_X X v_X$ \\
$\CS(\widetilde G)$ & a finite observation set exposing the typed reconstruction basis of $\widetilde G$ (Section~\ref{subsec:typed-local-control}) \\
$\mathcal B(\widetilde G)$ & the finite typed reconstruction basis of $\widetilde G$ (Definition~\ref{def:typed-basis}) \\
$\hat V(K)$ & the set of observed nonterminals $[x:u,v]$ extracted from factorizations $uxv\in K$ \\
$\hat G(K)$ & the canonical hypothesis grammar constructed from the finite sample $K$ (Section~\ref{subsec:learner}) \\
$\mathcal A_h$ & the learner defined by $\mathcal A_h(w_1,\dots,w_n)=\hat G(\{w_1,\dots,w_n\})$ \\
\bottomrule
\end{tabular}
\caption{Notation specific to typed reconstruction and learner construction.}
\label{tab:notation-typed}
\end{table}

This section proceeds in the order
$\text{typed local control}\to\text{finite typed basis}$. We first
introduce the typed refinement $\widetilde G$ of the target grammar and
show that there are only finitely many reachable productive typed
nonterminals and finitely many realized typed rule instances. We then
collect these into the finite typed reconstruction basis
$\mathcal B(\widetilde G)$ and show that this basis is exposed by a
finite observation set $\CS(\widetilde G)$.

\begin{definition}[start-separated binary normal form]\label{def:SSBNF}
A context-free grammar $G=(V,\Sigma,P,S_0)$ is in
\emph{start-separated binary normal form} (SSBNF) if the following
conditions hold: \textup{(i)} $S_0$ does not occur on the right-hand
side of any production; \textup{(ii)} every production with left-hand
side $S_0$ has the form $S_0\to A$ for some
$A\in V\setminus\{S_0\}$ or $S_0\to\lambda$; \textup{(iii)} every
production with left-hand side $A\in V\setminus\{S_0\}$ has the form
$A\to a$ for some $a\in\Sigma$ or $A\to BC$ for some
$B,C\in V\setminus\{S_0\}$.
\end{definition}

\begin{definition}[reduced grammar]\label{def:reduced}
A context-free grammar is \emph{reduced} if every nonterminal is
reachable from the start symbol and productive.
\end{definition}

Throughout the paper, we work with a reduced SSBNF grammar for the
target language. An equivalent reduced SSBNF grammar can be obtained
from any context-free grammar by the standard procedures of
start-symbol separation, elimination of useless symbols, and
binarization.

\subsection{Typed refinement}\label{subsec:typed-refinement}

Even when the underlying nonterminal symbol is the same, different
occurrences must be treated as different typed copies if the $h$-type of
the generated yield or the $h$-types of the surrounding left and right
contexts differ. In the fixed-$h$ setting, exact reconstruction depends
not only on what a nonterminal can generate, but also on where it can
appear with respect to the ambient recognizable congruence. Thus a
nonterminal must carry both an internal yield type and an external
left--right context type. Example~\ref{ex:typing-motivation} illustrates
this distinction concretely.

\begin{example}[why typing is needed]\label{ex:typing-motivation}
Let $G=(\{S,A,B\},\{a,b\},P,S)$ have productions
$S\to AB\mid BA$, $A\to a$, and $B\to b$. Then $L(G)=\{ab,ba\}$, but
the same nonterminal $A$ appears in two reachable sentential forms with
different context $h$-types. Thus, if we take the monoid
$M=(\mathbb{Z}/2\mathbb{Z},+,0)$ and define the homomorphism
$h:\{a,b\}^*\to M$ by $h(a)=0$ and $h(b)=1$, then the types distinguish
the contexts:
\[
\begin{array}{c|c|c}
\text{sentential form} & \text{context of $A$} &
(h(\text{left context}),\,h(\text{right context}))\\\hline
S\Rightarrow_G AB & (\lambda,b) & (0,1)\\
S\Rightarrow_G BA & (b,\lambda) & (1,0).
\end{array}
\]
The typing construction therefore separates these two occurrences of
$A$ into distinct typed copies, namely $A_0^{0,1}$ and $A_0^{1,0}$, so
that each typed nonterminal carries a globally consistent yield type and
outer context type.
\end{example}

We now define the typed refinement of the target grammar. Let
$L\in\Ccf{h}$, and let $G=(V,\Sigma,P,S_0)$ be any reduced SSBNF grammar
(see Definition~\ref{def:SSBNF}) such that $L(G)=L$. Write $e_M$ for the
identity element of $M$.

It is useful to distinguish two grammars. We first define the
\emph{full typed refinement} of $G$, denoted
\[
  \widetilde G^{\mathrm{full}}
  =
  (\widetilde V^{\mathrm{full}},\Sigma,
   \widetilde P^{\mathrm{full}},\widetilde S).
\]
Its start symbol is $\widetilde S$, and its non-start typed
nonterminals are the symbols $A_p^{m,n}$, where
$A\in V\setminus\{S_0\}$ and $m,n,p\in M$. The intended meaning is that
$A_p^{m,n}$ is a copy of $A$ whose generated yield has $h$-type $p$ and
whose surrounding left and right contexts have $h$-types $m$ and $n$,
respectively. The rules of $\widetilde G^{\mathrm{full}}$ are as
follows.
\begin{itemize}[nosep,leftmargin=2.2em]
\item If $S_0\to A\in P$, then include
$\widetilde S\to A_p^{e_M,e_M}$ for every $p\in M$.
\item If $S_0\to\lambda\in P$, then include
$\widetilde S\to\lambda$.
\item If $A\to a\in P$ and $h(a)=p$, then include
$A_p^{m,n}\to a$ for every $m,n\in M$.
\item If $A\to BC\in P$, then include
\[
  A_p^{m,n}\to B_q^{m,rn}\,C_r^{mq,n}
\]
for all $m,n,q,r\in M$ satisfying $qr=p$.
\end{itemize}

The superscripts in the binary rule express the propagation of outer
context types. If an occurrence of $A$ has left context type $m$ and
right context type $n$, and if the left child has yield type $q$ while
the right child has yield type $r$, then the left child sees right
context type $rn$, and the right child sees left context type $mq$.
The condition $qr=p$ records that the parent yield is the concatenation
of the two child yields.

The \emph{trimmed typed refinement} of $G$, denoted
\[
  \widetilde G=(\widetilde V,\Sigma,\widetilde P,\widetilde S),
\]
is obtained from $\widetilde G^{\mathrm{full}}$ by deleting all
unreachable symbols, all nonproductive symbols, and all rules incident
with such symbols. Unless explicitly stated otherwise,
$\widetilde G$ denotes this trimmed refinement.

The distinction between
$\widetilde G^{\mathrm{full}}$ and $\widetilde G$ is important. The full
refinement contains a typed copy $A_p^{m,n}$ for every
$A\in V\setminus\{S_0\}$ and every $m,n,p\in M$, regardless of whether
that copy can actually occur in a successful derivation. The trimmed
refinement keeps only the reachable and productive typed copies and
their realized rules. Thus derivations in
$\widetilde G^{\mathrm{full}}$ may mention typed copies that are later
removed, whereas all reconstruction data used below are taken only from
the trimmed grammar $\widetilde G$. In particular, the canonical yield
$\omega(X)$, the canonical context $\chi(X)$, the observation set
$\CS(\widetilde G)$, and the finite reconstruction basis
$\mathcal B(\widetilde G)$ are defined only for reachable and productive
symbols of $\widetilde G$. This convention avoids making any claim about
typed copies that do not exist after trimming.

\subsection{Canonical typed data}

Fix a total order on $\Sigma$, let $\le_{\mathrm{slex}}$ be the induced
shortlex order on $\Sigma^*$, and extend it lexicographically to
$(\Sigma^*)^2$. For each productive symbol
$\alpha\in(\Sigma\cup\widetilde V)^*$ of $\widetilde G$\footnote{That
is, a string of nonterminals from which one can eventually derive a
terminal string.}, let $\omega(\alpha)$ denote the shortlex-minimal
terminal string satisfying
$\alpha\Rightarrow_{\widetilde G}^*\omega(\alpha)$. In particular, if
$X\to YZ$, then $\omega(YZ)=\omega(Y)\omega(Z)$, and if $X\to a$, then
$\omega(a)=a$. Also, for each reachable non-start symbol
$X\in\widetilde V\setminus\{\widetilde S\}$, let
$\chi(X)=\langle u_X,v_X\rangle$ denote the shortlex-minimal pair
satisfying $\widetilde S\Rightarrow_{\widetilde G}^* u_X\,X\,v_X$.

The following typing invariants are immediate from the construction.

\begin{lemma}[full typing derivability]
\label{lem:typed-full}
In the full typed refinement $\widetilde G^{\mathrm{full}}$, if
$A\Rightarrow_G^* w\in\Sigma^+$ and $p=h(w)$, then for all
$m,n\in M$ we have
\[
  A_p^{m,n}\Rightarrow_{\widetilde G^{\mathrm{full}}}^* w .
\]
\end{lemma}

\begin{proof}
We argue by induction on the height of the derivation
$A\Rightarrow_G^* w$ in the untyped grammar. If the first rule is
$A\to a$, then $w=a$ and $p=h(a)$, so the rule
$A_p^{m,n}\to a$ belongs to $\widetilde G^{\mathrm{full}}$ for every
$m,n\in M$.

If the first rule is $A\to BC$ and $w=xy$ with
$B\Rightarrow_G^* x$ and $C\Rightarrow_G^* y$, put
$q:=h(x)$ and $r:=h(y)$. Then $qr=h(xy)=p$. By the induction hypothesis,
\[
  B_q^{m,rn}\Rightarrow_{\widetilde G^{\mathrm{full}}}^* x
  \quad\text{and}\quad
  C_r^{mq,n}\Rightarrow_{\widetilde G^{\mathrm{full}}}^* y .
\]
The full typed refinement contains the rule
\[
  A_p^{m,n}\to B_q^{m,rn}\,C_r^{mq,n},
\]
and hence $A_p^{m,n}\Rightarrow_{\widetilde G^{\mathrm{full}}}^*xy=w$.
\end{proof}

\begin{lemma}[typing invariants]
\label{lem:typed}
For the trimmed typed refinement $\widetilde G$, the following hold:
\textup{(i)} if $A_p^{m,n}\Rightarrow_{\widetilde G}^* w\in\Sigma^+$,
then $h(w)=p$;
\textup{(ii)} if $\widetilde S\Rightarrow_{\widetilde G}^* uA_p^{m,n}v$,
then $h(u)=m$ and $h(v)=n$;
\textup{(iii)} $L(\widetilde G)=L(G)=L$.
\end{lemma}

\begin{proof}
For \textup{(i)} and \textup{(ii)}, argue by induction on the height of
the derivation in the trimmed typed grammar. Since every rule of
$\widetilde G$ is a rule of $\widetilde G^{\mathrm{full}}$, the local
typing equations are exactly those specified in the full refinement.

For \textup{(i)}, the terminal case follows from rules
$A_p^{m,n}\to a$ with $p=h(a)$. In the binary case, if
\[
  A_p^{m,n}\to B_q^{m,rn}\,C_r^{mq,n}
\]
and the two children derive strings of $h$-types $q$ and $r$, then the
whole yield has $h$-type $qr=p$.

For \textup{(ii)}, the start case has outer context types
$(e_M,e_M)$ by construction. In the binary rule above, if the parent
occurs with left and right context types $(m,n)$, then the left child
has left context type $m$ and right context type $rn$, while the right
child has left context type $mq$ and right context type $n$, exactly as
prescribed by the rule. The claim follows by induction along the
derivation from $\widetilde S$.

For \textup{(iii)}, every derivation in $\widetilde G$ erases typed
annotations and gives a derivation in $G$, so
$L(\widetilde G)\subseteq L(G)$.

Conversely, let $w\in L(G)$. If $w=\lambda$, then $S_0\to\lambda$ is a
rule of $G$, so $\widetilde S\to\lambda$ is a rule of
$\widetilde G^{\mathrm{full}}$. Since this rule is reachable and
productive, it remains in the trimmed grammar $\widetilde G$, and hence
$\lambda\in L(\widetilde G)$.

Now suppose $w\neq\lambda$. Choose a derivation
$S_0\Rightarrow_G A\Rightarrow_G^* w$ and put $p:=h(w)$. In
$\widetilde G^{\mathrm{full}}$ the rule
$\widetilde S\to A_p^{e_M,e_M}$ exists, and by
Lemma~\ref{lem:typed-full} we have
$A_p^{e_M,e_M}\Rightarrow_{\widetilde G^{\mathrm{full}}}^* w$.
Therefore the symbol $A_p^{e_M,e_M}$ is reachable from $\widetilde S$
and productive in $\widetilde G^{\mathrm{full}}$. Moreover, the rules
appearing in one such successful derivation from $\widetilde S$ to $w$
are all incident only with reachable and productive symbols. Hence they
survive the trimming step, and the same derivation exists in
$\widetilde G$. Thus $w\in L(\widetilde G)$.
\end{proof}


\begin{lemma}\label{lem:typed-data}
If $X=A_p^{m,n}$ is reachable and productive in $\widetilde G$, then
$h(u_X)=m$, $h(\omega(X))=p$, and $h(v_X)=n$.
\end{lemma}

\begin{proof}
Since $\widetilde S\Rightarrow_{\widetilde G}^*u_XA_p^{m,n}v_X$,
Lemma~\ref{lem:typed}\textup{(ii)} gives $h(u_X)=m$ and $h(v_X)=n$.
Since $X\Rightarrow_{\widetilde G}^*\omega(X)$,
Lemma~\ref{lem:typed}\textup{(i)} gives $h(\omega(X))=p$.
\end{proof}

\subsection{Typed local control}\label{subsec:typed-local-control}

We extract the finite local typed information that will later be used for reconstruction. Define the characteristic sample by
\[
\CS(\widetilde G):=\{u_X\,\omega(\beta)v_X\mid X\to\beta\in\widetilde P\text{ and }X\in\mathrm{NT}(\widetilde G),\ X\neq\widetilde S\}\cup(\{\lambda\}\text{ if }\widetilde S\to\lambda\in\widetilde P),
\]
where $\mathrm{NT}(\widetilde G)$ denotes the set of reachable and productive typed nonterminals of $\widetilde G$.

\begin{lemma}[observation words are positive examples]
\label{lem:cs-contained-in-language}
We have $\CS(\widetilde G)\subseteq L(\widetilde G)=L$.
\end{lemma}

\begin{proof}
Let $z\in\CS(\widetilde G)$. If $z=\lambda$, then
$\widetilde S\to\lambda$ is a rule of $\widetilde G$, and hence
$\lambda\in L(\widetilde G)$.

Otherwise, by the definition of $\CS(\widetilde G)$, there is a
realizable non-start rule $X\to\beta$ of $\widetilde G$ such that
$z=u_X\omega(\beta)v_X$. Since $X$ is reachable, the definition of
$\chi(X)=\langle u_X,v_X\rangle$ gives
\[
  \widetilde S\Rightarrow_{\widetilde G}^* u_X X v_X .
\]
Since \(X\to\beta\) is a rule of the trimmed grammar and its right-hand
side is productive, we have
\[
  \beta\Rightarrow_{\widetilde G}^*\omega(\beta).
\]
Therefore
\[
  \widetilde S
  \Rightarrow_{\widetilde G}^* u_X X v_X
  \Rightarrow_{\widetilde G} u_X\beta v_X
  \Rightarrow_{\widetilde G}^* u_X\omega(\beta)v_X=z .
\]
Thus $z\in L(\widetilde G)$. By Lemma~\ref{lem:typed}\textup{(iii)},
$L(\widetilde G)=L$, so $z\in L$.
\end{proof}


\begin{lemma}\label{lem:typed-cs}
If $X$ is reachable and productive in $\widetilde G$, then $u_X\omega(X)v_X\in\CS(\widetilde G)$.
\end{lemma}

\begin{proof}
Choose a derivation $X\Rightarrow_{\widetilde G}^*\omega(X)$, and let $X\to\beta$ be the first rule used in this derivation. By definition, $\omega(X)=\omega(\beta)$. Hence $u_X\omega(X)v_X=u_X\omega(\beta)v_X\in\CS(\widetilde G)$.
\end{proof}

\begin{theorem}[Typed local control]\label{thm:typed-local-control}
$\CS(\widetilde G)$ is well defined and satisfies $|\CS(\widetilde G)|<\infty$. In other words, all local derivational information relevant to reconstruction is carried by finitely many typed local configurations.
\end{theorem}
\begin{proof}
The full typed refinement already has the finite nonterminal set
\[
  \{\widetilde S\}\cup
  \{A_p^{m,n}\mid A\in V\setminus\{S_0\},\ m,n,p\in M\}.
\]
The trimmed grammar $\widetilde G$ is obtained by deleting symbols and
rules from this finite grammar, so it is finite.

For each productive symbol $X$ of $\widetilde G$, the set
$\{w\in\Sigma^*\mid X\Rightarrow_{\widetilde G}^* w\}$ is nonempty, so
the shortlex least element $\omega(X)$ is well defined. Similarly, for
each reachable non-start symbol $X$ of $\widetilde G$, the set
\[
  \{\langle u,v\rangle\in(\Sigma^*)^2
    \mid \widetilde S\Rightarrow_{\widetilde G}^* uXv\}
\]
is nonempty, so the shortlex least pair
$\chi(X)=\langle u_X,v_X\rangle$ is well defined. Since $\widetilde G$
has only finitely many reachable and productive symbols and rules, the
set of observation words $u_X\omega(\beta)v_X$ appearing in the
definition of $\CS(\widetilde G)$ is finite. Therefore
$\CS(\widetilde G)$ is well defined and finite.
\end{proof}

\begin{remark}
The point of Theorem~\ref{thm:typed-local-control} is that, after typing, the information about the local behavioral conditions that the target grammar must satisfy for reconstruction is compressed into finitely many typed objects.
\end{remark}


\subsection{Finite typed reconstruction basis}

We now gather the finite typed local information of $\widetilde G$ into
a single reconstruction object.

\begin{definition}[finite typed reconstruction basis]\label{def:typed-basis}
Define the finite typed reconstruction basis of $\widetilde G$ by
\[
\mathcal B(\widetilde G):=\bigl(\mathrm{NT}(\widetilde G),\mathrm{Rule}(\widetilde G),\omega,\chi,\CS(\widetilde G)\bigr).
\]
Here, $\mathrm{NT}(\widetilde G)$ is the set of reachable and productive
typed nonterminals of $\widetilde G$, and
$\mathrm{Rule}(\widetilde G)$ is the set of rules whose symbols are all
reachable and productive typed nonterminals of $\widetilde G$ (we call
these the ``realizable rules'').
\end{definition}

\begin{theorem}[existence of a finite typed reconstruction basis]\label{thm:finite-typed-basis}
The basis $\mathcal B(\widetilde G)$ is well-defined and finite.
Moreover, each reachable and productive typed nonterminal is exhibited
by its anchor word $u_X\omega(X)v_X\in\CS(\widetilde G)$, and each
realizable rule $X\to\beta$ is exhibited by its canonical observation
word $u_X\omega(\beta)v_X\in\CS(\widetilde G)$.
\end{theorem}

\begin{proof}
Finiteness follows from Theorem~\ref{thm:typed-local-control}. By
Lemma~\ref{lem:typed-cs}, each reachable and productive typed
nonterminal $X$ is exhibited by its anchor word
$u_X\omega(X)v_X$. Also, by the definition of $\CS(\widetilde G)$, each
realized non-start rule $X\to\beta$ is exhibited by the corresponding
canonical observation word $u_X\omega(\beta)v_X$.
\end{proof}

\begin{remark}
In this notation, the characteristic sample $\CS(\widetilde G)$ is the
terminal observation set that exhibits the finite typed reconstruction
basis $\mathcal B(\widetilde G)$; that is, it serves as an observation
window through which the hidden typed structure is made visible. Thus
the characteristic sample is by no means merely an auxiliary set added
later only for the convergence proof. Rather, it is precisely the
working space, or finite observational interface, on which the fixed-$h$
reconstruction theory operates.
\end{remark}
